\documentclass[10pt,a4paper]{dinbrief}
\usepackage[utf8]{inputenc}
\usepackage{ngerman}
\usepackage[
top    = 2.75cm,
bottom = 2.50cm,
left   = 3.00cm,
right  = 2.50cm]{geometry}
\usepackage{eurosym}
\usepackage{fancyhdr}
\usepackage{pdfpages}
\address{% 
	Dr. Tobias Griebe\\
	Gartenstra{\ss}e 29\\
	15366 Neuenhagen 
}
\backaddress{Dr. T. Griebe, Gartenstr. 29, 15366 Neuenhagen}
\signature{Dr. Tobias Griebe}
\place{Neuenhagen}
\nowindowrules \nobackaddressrule \nowindowtics %\normaladdress
\date{\today}
\pagestyle{empty}
\begin{document} 
\begin{letter}  
{%
	IT-Dienstleistungszentrum Berlin\\
	Berliner Straße 112-115\\
	10713 Berlin
}
\subject{\bf Bewerbung Abteilungsleiter Kennziffer 70/2019}
\opening{Sehr geehrte Damen und Herren,}
mit großem Interesse bin ich auf Ihre Stellenausschreibung Abteilungsleiter beim ITDZ Berlin mit Kennziffer 70/2019 gestoßen. 

Seit beinahe 10 Jahren bearbeite ich IT-Projekte im Themenkomplex Digitalisierung in meinen Tätigkeiten als Berater und Führungskraft im Kernbereich Digitalisierung bei der adesso AG sowie zuvor in meiner akademischen Tätigkeit als Mitarbeiter an den Universitäten Leipzig und Duisburg-Essen. 

Die im BEGS beschriebene Weiterentwicklungen elektronischer Behördengänge und der Einsatz moderner Informationstechnik in der Berliner Verwaltung fällt damit in einen Themenkomplex, den ich als Handlungsfeld von zukunftsweisender Bedeutung mit großem Einfluss auf die gesellschaftliche und politische Entwicklung der nächsten Jahrzehnte betrachte.

Als Abteilungsleiter bei der adesso AG befasse ich mich mit Projekten im gesamten Handlungsfeld interdisziplinärer Digitalisierungsprojekte unserer Wirtschaftspartner sowie mit der Begleitung der Strategieentwicklung, des Projektmanagements und der Koordination von IT-Projekten in strategischen und operativen Rollen sowie mit der Steuerung interner und externer Entwicklungsdienstleister. 

Durch meine Tätigkeit bei der adesso AG und zuvor als Leiter der Arbeitsgruppe Mobile Anwendungen und Prozesse an der Universität Duisburg-Essen blicke ich auf mehrjährige Erfahrung in der fachlichen und disziplinarischen Mitarbeiterführung zurück. Insbesondere verantworte ich in meiner aktuellen Position den Aufbau und die Führung eines Teams von Mitarbeitern, innerhalb meiner eigenen Organisationseinheit und übergreifend im Projektgeschäft. Hier übernehme ich alle Aktivitäten zur Entwicklung von Mitarbeitern von Recruiting, Qualifikation, Laufbahnstufenentwicklung bis hin zur Retention, Identifikation von Leistung und Potenzial meiner Mitarbeiter, Ermittlung von Entwicklungsbedarf sowie dem zielgerichteten Einsatz passender Personalentwicklungsinstrumente.

Die Analyse von Bedarfsentwicklungen in unterschiedlichen Themenfeldern, die Evaluation von Anbietermärkten und Branchenlösungen sowie die Kooperation mit Dienstleistern und Softwareherstellern gehören zu meiner täglichen Praxis bei der Umsetzung von Konzepten und Lösungen. Hierzu gehört ebenfalls die fachliche und strategische Beratung auf Vorstandsebene und die Repräsentation meines Unternehmens bei Kunden und Partnern. Der sichere Umgang mit klassischen und agilen Vorgehensmodellen gehört dabei ebenso zu meinen Kompetenzen wie der Umgang mit State-of-the-Art Werkzeugen zur aktiven und passiven Begleitung/Controlling von IT-Projekten. Im Rahmen der von mir verantworteten Projekten setze ich ebenfalls meine fundierten Kenntnisse in den Bereichen Hardware, Software, Programmiersprachen, Architekturen und Cloud-basierte Backend-Systeme regelmäßig ein.

Nach meinen Erfahrungen in der Beratung bin ich nun auf der Suche nach einem neuen Engagement. In meinen nächsten Karriereschritt möchte ich meine Tätigkeiten in einer gesunden Mischung aus disziplinarischer Mitarbeiterführung und spannender Mitgestaltung von IT-Lösungen im Kontext der Digitalisierung am Standort Berlin fortsetzen. Besonders interessant sind für mich Projektumgebungen mit erkennbarem praktischem Nutzen für zukünftige Anwender.

Vor diesem Hintergrund bin ich davon überzeugt, durch meinen ausgeprägten technischen und akademischen Hintergrund, meine eigenständige und qualitätsorientierte Arbeitsweise, meine Kommunikationsbereitschaft, Organisationsfähigkeit und Dienstleistungsorientierung sowie meine Fähigkeit Mitarbeiter anzuleiten, zu entwickeln und zu motivieren zu Ihrem Erfolg beizutragen.

Ich bin ab Oktober 2019 verfügbar (verhandelbar).

Über eine Einladung zu einem persönlichen Gespräch freue ich mich sehr.


\closing{Mit freundlichen Gr"u"sen,}
\encl{Lebenslauf, Promotionsurkunde, Diplomurkunde, Zertifikat Licence to Lead F"uhrungskr"afteausbildung, Zertifikat ISTQB Certified Tester, Zertifikat Scrum.org Professional Scrum Product Owner, Arbeitszeugnisse}
\end{letter}
\end{document}